\documentclass[15pt]{scrartcl}
\usepackage[sexy]{evan}
\newcommand{\R}{\mathbb{R}}
\newcommand{\N}{\mathbb{N}}
\newcommand{\Q}{\mathbb{Q}}
\usepackage{relsize}
\begin{document}
\title{H104 Midterm Review}
\date{\today}
\author{Jiawen Huang}
\maketitle

\section{Character of a field}
\begin{enumerate}
    \item Morphism between fields must be injective.
    \item Character of a field must be prime or $0$.
    \item If the character is $0$, the field contains a copy of $\Q$.
    \item If the character is a prime $p$, the field contains a copy of $\mathbb{F}_p$. 
    \item If there is morphism from $F$ to $K$, then they must have the same character.
    \begin{remark}
    Fields with different characters are basically unrelated.
    \end{remark}
\end{enumerate}
\section{Field Extensions}
\begin{enumerate}
    \item Every field is an extension of its prime field or $\Q$.
    \item $k\subseteq F$. $F$ is a vector space over $k$ and moreover a $k$-algebra.
    \item The degree of extension $[F:k]$ is the dimension of $F$ as a vector space over $k$.
    \item If $k\subseteq F \subseteq K$, then $[K:k] = [K:F][F:k]$. 
\end{enumerate}
\section{Simple Extensions}
\begin{enumerate}
    \item An extension $F/k$ is simple if $F = k(\alpha)$ for some $\alpha \in F$. $k(\alpha)$ is the smallest subfield of $F$ containing $k$ and $\alpha$.
    \item Examples of simple extensions: $\Q(\sqrt{2})$, $\mathbb{F}_p(t)$. Is $\Q(\sqrt{2}, \sqrt{3})$ simple? Yes, $\Q(\sqrt{2}, \sqrt{3}) = \Q(\sqrt{2} + \sqrt{3})$ (look at minimal polynomial of $\sqrt{2} + \sqrt{3}$ over $\Q$ and the degree of $\Q(\sqrt{2}, \sqrt{3})$ over $\Q$).
    \item Surprisingly, \textbf{every finite extension of character 0 is simple !}. In fact, it's very hard to find an example of non-simple extension.
    \item So what does $k(\alpha)$ look like? Consider the evaluation morphism $f: k[x] \to F$ defined by $f(p(x)) = p(\alpha)$. The kernel of this morphism is an ideal in $k[x]$, and since $k[x]$ is a principal ideal domain, the kernel is generated by a single polynomial, called the minimal polynomial of $\alpha$ over $k$, and we say $\alpha$ is algebraic over $k$. If the minimal polynomial is $0$, then $k(\alpha)$ is isomorphic to the field of rational functions $k(x)$. 
    \item The degree of the minimal polynomial of $\alpha$ over $k$ is equal to $[k(\alpha):k]$, why?
\end{enumerate}
\section{Algebraic Number}
If $F=k(\alpha_1, \alpha_2, \ldots, \alpha_n)$ (finitely generated), then the following are equivalent:
\begin{enumerate}
    \item $[F:k]$ is finite.
    \item $F$ is made of entirely algebraic elements over $k$.
    \item Each $\alpha_i$ is algebraic over $k$.
\end{enumerate}
Key of proof: $k \subseteq k(\alpha_1) \subseteq k(\alpha_1,\alpha_2) \subseteq \cdots \subseteq k(\alpha_1,\ldots,\alpha_n) = F$
\begin{remark}
    This is very powerful. 
\begin{itemize}
    \item If $a$ and $b$ are algebraic over $k$, then $a+b$, $a-b$, $ab$, and $a/b$ (if $b \neq 0$) are also algebraic over $k$.
    \item Let $E$ be the set of algebraic numbers over $k\subseteq F$, then $E$ is a field. Furthermore, $E$ is an extension of $k$ called the algebraic closure of $k$ in $F$ (why?).
\end{itemize}

\end{remark}
\section{splitting field}
\begin{enumerate}
    \item Unique up to isomorphism. We skip the proof.
    \item Splitting field must be a finite extension.
    \item Example: splitting field of $x^4+1$ over $\Q$ is $\Q(\sqrt{2}, i)$.
    \item Example: splitting field of $x^4+2$ over $\Q$ is $\Q(\sqrt[4]{2}, i)$.
\end{enumerate}


\end{document}